\documentclass[10pt]{article}

\usepackage[left=0.85in, right=0.85in, top=0.7in, bottom=0.7in]{geometry}
\usepackage{amsmath}
\usepackage{booktabs}
\usepackage{titlesec}
\usepackage{parskip}
\usepackage{microtype}
\usepackage{enumitem}
\usepackage{times}

\setlength{\parskip}{2.5pt}
\setlength{\parindent}{0pt}
\setlength{\columnsep}{0.23in}

\titleformat{\section}{\normalsize\bfseries}{}{0em}{}
\titlespacing{\section}{0pt}{6pt}{2pt}

\title{\textbf{CRAFT: A Self-Building Agent Framework for Telecom Configuration}}
\author{Anonymous Authors}
\date{}

\begin{document}
\maketitle

% -------- ABSTRACT --------
\onecolumn
\section*{Abstract}

Telecom network configuration is manual, repetitive, and operator-specific - each new operator requires rebuilding logic from scratch. \textbf{CRAFT} (Configuration RAN Auto-generation Framework for Telecom) introduces six AI agents that learn rules from past deployments, reason with 3GPP standards, discover templates, test rules, build missing tools, and continuously improve. Evaluated on Vodafone and Telus deployments, CRAFT achieves \textbf{95.2\% field accuracy}, reduces generation time from 2-3 days to 4 minutes, and reaches \textbf{78.6\% rule coverage on a new operator with zero prior exposure}.

\twocolumn

\section{Introduction}

Telecom configuration remains a per-operator engineering task. Engineers manually map CIQ, LLD, and rule sheets into scripts, repeating the process for every operator. Existing LLM approaches reduce effort but do not learn from deployments or generalize across operators.

CRAFT treats configuration as a knowledge extraction problem. It leverages past deployments, domain standards, and cross-operator patterns through a coordinated multi-agent system.

\section{Framework Overview}

CRAFT consists of five stages and a feedback loop:

\textbf{ISA (Learning)} extracts rules from historical deployments. Stable values become defaults; variable fields are traced to CIQ or LLD sources.

\textbf{DGRA (Reasoning)} resolves complex rules using retrieval over 3GPP standards and cross-operator knowledge. It handles multi-source dependencies and transfers rules across operators.

\textbf{TDA (Template Discovery)} identifies required output templates based on configuration context.

\textbf{MTA (Validation)} stress-tests rules using boundary and conflicting inputs.

\textbf{PVS (Tool Synthesis)} builds missing computation tools via code generation and sandbox validation.

\textbf{FLA (Feedback)} captures corrections and updates rules globally.

Together, ISA and DGRA achieve over \textbf{90\% rule coverage} before template execution.

\section{Results}

Evaluated on two operators:

\textbf{Vodafone:} CIQ + LLD + 81 rules.

\textbf{Telus:} heterogeneous deployment files with no prior exposure.

\vspace{2pt}
\begin{center}
\small
\renewcommand{\arraystretch}{1.2}
\begin{tabular}{lcc}
\toprule
Metric & Vodafone & Telus \\
\midrule
Rule coverage & 91.4\% & 78.6\% \\
Accuracy & 95.2\% & - \\
Tool success & 5/6 & - \\
Time & 4 min & 4 min \\
Manual baseline & 2-3 days & 2-3 days \\
\bottomrule
\end{tabular}
\end{center}
\vspace{2pt}

Despite zero prior exposure, CRAFT achieves 78.6\% rule coverage on Telus through cross-operator learning.

\section{Impact}

\textbf{Faster onboarding:} reduced from 6-12 months to weeks.

\textbf{Error reduction:} below 0.5\% via validation and confidence gating.

\textbf{Knowledge compounding:} each deployment improves future performance.

\textbf{AI-native systems:} enables autonomous configuration workflows.

\section{Conclusion}

CRAFT transforms telecom configuration from a manual coding task into a self-improving knowledge system. By combining learning, reasoning, validation, and feedback, it achieves rapid, accurate, and scalable configuration across operators.

\begin{thebibliography}{3}
\small
\bibitem{lira2024}
O.~G. Lira et al., ``LLMs for network configuration,'' arXiv, 2024.

\bibitem{provvedi2026}
C.~Provvedi et al., ``Intent-LLM,'' IEEE TCCN, 2026.

\bibitem{wang2026}
J.~Wang et al., ``LLM-driven configuration,'' IEEE TNSM, 2026.
\end{thebibliography}

\end{document}
